\chapter{多変数の微分積分}
\label{ch:tb-multivariable}

\section{連続性、偏微分、\texorpdfstring{$C^1$ 級}{C1 級}}

$C^1$ 級は1階偏導関数が存在して連続、$C^\infty$ 級は全階の偏導関数が連続、という意味である。
$C^1$ 級なら連続だが、偏導関数が存在するだけでは連続とは限らない。

原点で場合分けされた関数は、次の順に処理する。
\begin{enumerate}
  \item 連続性：$x=r\cos\theta$, $y=r\sin\theta$ とし、
        $\abs{f(x,y)-f(0,0)}\le Cr^\alpha$（$\alpha>0$）へ落とす。
  \item 原点での偏微分：周囲の式を形式微分せず、軸上の差商
        \[
        f_x(0,0)=\lim_{h\to0}\frac{f(h,0)-f(0,0)}h,
        \quad
        f_y(0,0)=\lim_{h\to0}\frac{f(0,h)-f(0,0)}h
        \]
        を計算する。
  \item $C^1$ 級でないこと：周囲で偏導関数を計算し、原点へ近づける経路を一つ選んで
        不連続を示す。
\end{enumerate}
評価では $\abs{\sin u}\le\min\{\abs u,1\}$ を、原点付近と無限遠で使い分ける。

\begin{theorem}{曲線に沿う連鎖律}{tb-multivariable-chain-rule}
$f$ が $C^1$ 級で $x(t),y(t)$ が微分可能なら
\[
\frac{\d}{\d t}f(x(t),y(t))
=f_x(x(t),y(t))x'(t)+f_y(x(t),y(t))y'(t).
\]
\end{theorem}

\begin{proposition}{原点で消える滑らかな関数の分解}{tb-smooth-factorization}
$f\in C^\infty(\R^2)$, $f(0,0)=0$ なら
\[
f(x,y)=xg(x,y)+yh(x,y)
\]
となる $g,h\in C^\infty(\R^2)$ が存在する。さらに
\[
f(x,y)=xG(x,y)+yH(y)
\]
という形にもできる。
\end{proposition}

\begin{proof}
$F(t)=f(tx,ty)$ とおく。基本定理と連鎖律から
\[
f(x,y)=\int_0^1F'(t)\,\d t
=x\int_0^1f_x(tx,ty)\,\d t
+y\int_0^1f_y(tx,ty)\,\d t.
\]
これが最初の表示である。二つ目は
\[
f(x,y)=\bigl(f(x,y)-f(0,y)\bigr)+f(0,y)
=x\int_0^1 f_x(tx,y)\,\d t
+y\int_0^1f_y(0,ty)\,\d t
\]
とすればよい。被積分関数が滑らかなので、得られた関数も滑らかである。
\end{proof}

\section{積分記号下の微分}

積分区間が動くか、中身だけが動くかを最初に分ける。

\begin{theorem}{Leibniz 則}{tb-leibniz-rule}
$f,f_x$ が必要な長方形上で連続なら
\[
\frac{\d}{\d x}\int_c^d f(x,y)\,\d y
=\int_c^d f_x(x,y)\,\d y.
\]
さらに $\phi,\psi$ が $C^1$ 級なら
\[
\frac{\d}{\d x}\int_{\phi(x)}^{\psi(x)}f(x,y)\,\d y
=f(x,\psi(x))\psi'(x)-f(x,\phi(x))\phi'(x)
+\int_{\phi(x)}^{\psi(x)}f_x(x,y)\,\d y.
\]
\end{theorem}

右辺は順に「上端」「下端」「中身」の変化である。
証明は $Q(x,v)=\int_c^v f(x,y)\,\d y$ と置き、
$Q(x,\psi(x))-Q(x,\phi(x))$ に二変数の連鎖律を使う。

例えば
\[
f(x)=\int_a^x\d s\int_a^s u(s,t)\,\d t
\]
では、外側の基本定理だけで
$f'(x)=\int_a^xu(x,t)\,\d t$ となる。

\begin{theorem}{パラメータ積分の微分}{tb-differentiate-integral}
$t$ がコンパクト区間を動き、$f(t,x)$ と $f_t(t,x)$ が連続で、
$\abs{f_t(t,x)}\le g(x)$ となる可積分関数 $g$ があれば
\[
\frac{\d}{\d t}\int f(t,x)\,\d x=\int f_t(t,x)\,\d x.
\]
広義積分では、交換を正当化する優関数 $g$ を答案に書く。
\end{theorem}

答案では「形式微分したら簡単になった」で終えない。
$t$ の近傍を固定し、そこで共通に使える可積分な $g(x)$ を一つ書く。

\begin{example}{$(x^t-1)/\log x$}{tb-log-quotient-parameter}
$0<x<1$, $t\ge0$ では
\[
\frac{x^t-1}{\log x}=\int_0^t x^s\,\d s.
\]
この表示から $0\le (x^t-1)/\log x\le t$、$x\to1$ での極限が $t$ と分かる。
さらに $t$ で微分すれば $x^t$ なので
\[
\frac{\d}{\d t}\int_0^1\frac{x^t-1}{\log x}\,\d x
=\int_0^1x^t\,\d x=\frac1{t+1}.
\]
積分表示は端点の特異性、評価、微分を同時に処理する。
\end{example}

\section{重積分と変数変換}

被積分関数より先に領域を見る。直交座標のまま上下端が一行で書けるなら累次積分、
円・和・比が境界に現れるなら変数変換を選ぶ。

\begin{theorem}{Fubini--Tonelli}{tb-fubini-tonelli}
$f\ge0$ なら Tonelli の定理により累次積分の順序を交換できる。
符号を持つ $f$ では $\iint\abs f<\infty$ なら Fubini の定理により交換できる。
有界閉領域上の連続関数では後者の条件が自動的に成り立つ。
\end{theorem}

累次積分へ直す前に領域を図示し、例えば
\[
D=\{(x,y)\mid a\le x\le b,\ \phi(x)\le y\le\psi(x)\}
\]
の形へ書く。境界の上下関係が変わる点で積分区間を分ける。

\begin{theorem}{二変数の変数変換}{tb-change-of-variables}
$\Phi:(u,v)\mapsto(x,y)$ が領域間の $C^1$ 級全単射で、内部で Jacobian が消えないなら
\[
\iint_D f(x,y)\,\d x\d y
=\iint_{\Phi^{-1}(D)}f(\Phi(u,v))
\left|\frac{\partial(x,y)}{\partial(u,v)}\right|\d u\d v.
\]
掛けるのは
\[
\left|\frac{\partial(x,y)}{\partial(u,v)}\right|
=\left|\det
\begin{pmatrix}x_u&x_v\\y_u&y_v\end{pmatrix}\right|
\]
であり、絶対値を忘れない。
\end{theorem}

変数変換は次の順に行う。
\begin{enumerate}
  \item 逆変換を求める。
  \item 領域の全ての不等式を新変数へ移し、新領域を確定する。
  \item Jacobian の絶対値を計算する。
  \item 被積分関数と面積要素を両方置き換える。
\end{enumerate}

\begin{proposition}{頻出する変換}{tb-standard-transformations}
\[
\begin{array}{c|c|c}
\text{目的}&(x,y)&\left|\partial(x,y)/\partial(u,v)\right|\\ \hline
\text{極座標}&(r\cos\theta,r\sin\theta)&r\\
x+y,\ x-y&( (u+v)/2,(u-v)/2)&1/2\\
\text{合計と割合}&(st,s(1-t))&s\\
\text{平方した極座標}&((r\cos\theta)^2,(r\sin\theta)^2)&4r^3\sin\theta\cos\theta
\end{array}
\]
最後の二行では通常 $s>0$, $0<t<1$ および $r>0$, $0<\theta<\pi/2$ を使う。
\end{proposition}

極座標で $0\le x\le1$ は $0\le r\cos\theta\le1$ となる。
$1\le r\le2$ と同時に課されるなら $r\le\min\{2,\sec\theta\}$ であり、
$\sec\theta=2$ となる $\theta=\pi/3$ で領域を分ける。

\begin{proposition}{対称性}{tb-symmetry-vanishing}
領域 $D$ が変換 $T$ で不変、Jacobian の絶対値が1で、
$f(Tx)=-f(x)$ なら $\int_Df=0$ である。
特に原点対称な領域上では、奇関数の積分は0になる。
\end{proposition}

\begin{proposition}{Gaussian 積分}{tb-gaussian-integral}
\[
\int_0^\infty e^{-x^2}\,\d x=\frac{\sqrt\pi}{2}.
\]
\end{proposition}

\begin{proof}
$I_R=\int_0^Re^{-x^2}\,\d x$ とし、$S_R=[0,R]^2$ とおくと Fubini の定理により
\[
I_R^2=\iint_{S_R}e^{-(x^2+y^2)}\,\d x\d y.
\]
第1象限の半径 $a$ の四分円を $D_a$ とすれば
$D_R\subset S_R\subset D_{\sqrt2R}$ である。極座標により
\[
\iint_{D_a}e^{-(x^2+y^2)}\,\d x\d y
=\int_0^{\pi/2}\!\int_0^a e^{-r^2}r\,\d r\d\theta
=\frac\pi4(1-e^{-a^2}).
\]
両側から挟んで $R\to\infty$ とすると $I_R^2\to\pi/4$。
$I_R\ge0$ なので結論を得る。
\end{proof}

\section{Gamma・Beta 関数}

\begin{definition}{Gamma・Beta 関数}{tb-gamma-beta}
$p,q>0$ に対し
\[
\Gamma(p)=\int_0^\infty x^{p-1}e^{-x}\,\d x,
\qquad
B(p,q)=\int_0^1t^{p-1}(1-t)^{q-1}\,\d t.
\]
\end{definition}

\begin{theorem}{Beta--Gamma 関係式}{tb-beta-gamma-relation}
\[
B(p,q)=\frac{\Gamma(p)\Gamma(q)}{\Gamma(p+q)}.
\]
\end{theorem}

\begin{proof}
積を第1象限の積分
\[
\Gamma(p)\Gamma(q)=\iint x^{p-1}y^{q-1}e^{-(x+y)}\,\d x\d y
\]
にする。$x=st$, $y=s(1-t)$ を使えば
\[
\int_0^\infty s^{p+q-1}e^{-s}\,\d s
\int_0^1t^{p-1}(1-t)^{q-1}\,\d t
=\Gamma(p+q)B(p,q).
\]
別解では $x=(r\cos\theta)^2$, $y=(r\sin\theta)^2$ とし、
Jacobian $4r^3\sin\theta\cos\theta$ を掛ける。$r$ と $\theta$ の積分に分離した後、
$s=r^2$, $t=\cos^2\theta$ と置けば同じ二因子を得る。
\end{proof}

\section{円板を広げる極限}

2021年度第1期では、非負関数を円板 $B(R)$ 上で積分する。
$B(R)$ は $R$ とともに増大するので、単調収束定理から
\[
\lim_{R\to\infty}\iint_{B(R)}F
=\iint_{\R^2}F
\]
となる。右辺が有限であることは原点と無限遠を分けて評価する。
問題の関数では、極座標で
\[
\abs f\le r\quad(0<r\le1),\qquad
\abs f\le r^{-2}\quad(r\ge1)
\]
なら、$p>1$ に対し
\[
\iint\abs f^p
\le2\pi\left(\int_0^1r^{p+1}\,\d r+\int_1^\infty r^{1-2p}\,\d r\right)<\infty.
\]
原点と無限遠で別の評価を使う典型である。
